\documentclass{article}
\usepackage[T1]{fontenc}
\usepackage{lmodern}
\usepackage{graphicx}
\usepackage{amsmath,amssymb}
\usepackage[a4paper,margin=2cm]{geometry}
\usepackage[absolute,overlay]{textpos}
\setlength{\TPHorizModule}{1mm}
\setlength{\TPVertModule}{1mm}
\usepackage[indonesian]{babel}
\usepackage{hyperref}
\setlength{\emergencystretch}{2em}
% unit: Halaman 53

\begin{document}

% === Halaman 53 (python/native) ===

53

j = 0;

 for (i=0; i<256; i++) \{

     j = (j + S[i] + K[i \% K.length()]) \% 256;

     swap(S[i], S[j]); // Pertukarkan nilai S[i] dan S[j] 

 \}

 i = 0;   j = 0;

while (!feof(Fin)) \{

     c = fgetc(Fin); 

     i = (i + 1) \% 256;

     j = (j + S[i]) \% 256;

     swap(S[i], S[j]);    // Pertukarkan nilai S[i] dan S[j] 

     t = (S[i] + S[j]) \% 256;

     u = S[t];   // keystream 

     p = u \textasciicircum{} c;       

     fputc(p, Fout);     

 \}  

 fclose(Fin);  fclose(Fout);

\}

\end{document}
